\section{Fluid Dynamics and Equilibrium}\label{sec:fluid}
Unlike traditional queueing systems, LLM inference features dynamically growing memory: KV cache expands as prompts progress through decode stages. This makes predicting overflow difficult; when total memory exceeds capacity $C$, evictions waste partial computation, degrading throughput and latency.

We develop a fluid model characterizing equilibrium---a balanced state where arrivals match completions. This equilibrium serves two purposes: it quantifies the maximum sustainable throughput $\throughput^*$ under memory constraint $C$, and it reveals the load-balancing principle guiding our algorithm (Section~\ref{sec:wait}). The principle is simple: maintain the system near equilibrium to prevent evictions. We first analyze single-type systems, then extend to multiple heterogeneous types.

\subsection{Single-Type Fluid Model}
For a single type \(j\), each prompt uses \(l_j\) tokens of memory at prefill stage ($s=0$) and \(l_j + s\) tokens at decode stage \(s = 1, \dots, l_j'\). At equilibrium, the system maintains \(n_j^*\) active prompts evenly distributed across $(l_j'+1)$ stages. This balanced distribution minimizes overflow risk: if prompts accumulate at later stages (higher memory), total usage may exceed $C$, forcing evictions.

The system becomes unstable when arrivals exceed the inverse of total processing time $d_1 (l_j' + 1)(l_j + l_j'/2)$.

\begin{proposition}\label{prop:large_rate}
When the condition $\lambda_j d_1 (l_j' + 1) \left( l_j + \frac{l_j'}{2} \right) \geq 1$ is satisfied, the system becomes unstable in the sense that the expected average latency under any scheduling policy with deterministic arrival rate \( \lambda_j \) grows linearly with time. Formally,
$$
\mathbb{E}[\textbf{Latency}^{(T,\pi)}] = \Omega(T) \quad \text{as} \quad T \to \infty.
$$
\end{proposition}

Proof in Appendix~\ref{appendix:proof_large_rate}. We assume $\lambda_j d_1 (l_j' + 1) \left( l_j + \frac{l_j'}{2} \right) <1$ henceforth, ensuring equilibrium exists.

Total memory at equilibrium is:
\begin{equation}
    M_j^* = \frac{n_j^*}{l_j^{\prime}+1}\sum_{s=0}^{l_j^{\prime}}(l_j+s) = n_j^*\prn{l_j+\frac{l_j'}{2}},
\end{equation}
where $(l_j + l_j'/2)$ is average memory per prompt. At equilibrium, arrivals equal completions:
$$ \underset{\text{Time per iteration}}{\underbrace{(d_0+d_1M_j^{*} )}}\lambda_j =\left(d_0+d_1 n_j^*\left( l_j +\frac{l_j^{\prime}}{2} \right) \right)\lambda_j= \frac{n_j^*}{l_j^{\prime}+1}. $$

Solving for $n_j^*$ yields:
$$n_j^* = \frac{d_0\lambda_j(l_j^{\prime}+1)}{1-d_1\lambda_j(l_j^{\prime}+1)(l_j+\frac{l_j^{\prime}}{2})}.$$

Equilibrium memory is:
\begin{equation}
    \label{eq:memory_single}
M_j^* = n_j^*\left( l_j+\tfrac{l_j^{\prime}}{2} \right)
=  \frac{d_0\lambda_j(l_j^{\prime}+1)\left( l_j+\tfrac{l_j^{\prime}}{2} \right) }{1-d_1\lambda_j(l_j^{\prime}+1)(l_j+\tfrac{l_j^{\prime}}{2})}.
\end{equation}
% \begin{comment}
% When the available memory is bigger than the equilibrium memory $M_j^*$, we can form a batch with a bigger  batch size. For example, each stage has $\frac{n_j^*}{l_j^\prime+1}+1$ prompts, and the memory $M_j^{+}$ is that
% $$M_j^{+} = \left(\frac{n_j^*}{l_j^\prime+1}+1\right)\sum_{k=0}^{l_j^{\prime}}(l_j+k)=(n_j^*+l_j^\prime+1)\left( l_j + \frac{l_j^\prime}{2} \right) = M_j^* + (l_j^\prime+1)\left( l_j + \frac{l_j^\prime}{2} \right)  $$
% So the arrival and completion ratio is that 
% $$ 
% \frac{\text{Arrival}}{\text{Completion}} 
% = \frac{\left(d_0+d_1M_j^*+d_1 (l_j^\prime+1)\left( l_j + \frac{l_j^\prime}{2} \right)\right)\lambda_j}{\frac{n_j^*}{l_j^\prime+1}+1} 
% = \frac{n_j^*+d_1 \lambda_j(l_j^\prime+1)^2\left( l_j + \frac{l_j^\prime}{2} \right)}{n_j^*+l_j^\prime+1} < 1 $$
% Where the last inequality comes from that $d_1 \lambda_j(l_j^\prime+1)\left( l_j + \frac{l_j^\prime}{2} \right)<1$. Similarly, when the available memory is smaller than equilibrium memory $M_j^*$, we have to form a batch with a smaller batch size. For example, each stage has $\frac{n_j^*+1}{l_j^\prime} -1$ prompts, and the memory $M_j^{-}$ is that
% $$
% M_j^{-}
% = \left( \frac{n_j^*}{l_j^\prime+1}-1 \right) \sum_{k=0}^{l_j^\prime}(l_j+k) = M_j^* - (l_j^\prime+1)\left(l_j + \frac{l_j^\prime}{2} \right)
% $$
% So the arrival and completion ratio is that 
% $$
% \frac{\text{Arrival}}{\text{Completion}} 
% = \frac{n_j^* - d_1\lambda_j(l_j^\prime+1)^2\left( l_j + \frac{l_j^\prime}{2} \right) }{ n_j^* - (l_j^\prime+1)  }>1
% $$
% \end{comment}
Equilibrium throughput is:
\begin{equation}\label{eq:throughput_fluid}
\begin{aligned}
\throughput_j^* &= \frac{n_j^* \cdot l_j'}{(l_j'+1)(d_0+d_1n_j^*\prn{l_j+\tfrac{l_j'}{2}})} \\
&= \lambda_j l_j',
\end{aligned}
\end{equation}
where the equality follows from the equilibrium condition.

% \gan{I should add more explanation. Why we introduce this class.}
% Moreover, we define a class of scheduling policies \(\Pi\) that includes non-predictive policies without \textit{preemption}. Here, \textit{non-predictive} means the policy uses only the realized sample path up to the current time, without knowing future information. The term \textit{preemption} describes the action where the policy clears the memory used by an ongoing prompt request and re-queues it to the front of the waiting prompts, usually because of limited memory. Policies in \(\Pi\) avoid preemption because it causes problems. For instance, preemption interrupts the ongoing work, requiring the system to restart or reload data for that prompt, which takes extra time and uses more computing resources. This delay in finishing the preempted prompt also lowers the system’s overall efficiency by repeating tasks unnecessarily. As an example, a non-predictive policy that reserves memory of \(l_{\max} + l_{\max}'\) for each prompt—where \(l_{\max} + l_{\max}'\) is the maximum memory one prompt can need—avoids preemption by always having enough memory, making it a policy in \(\Pi\) (though in this way, it may utilize the memory inefficiently).
% \gan{cite some example of scheduling policies without Preemption.}

% \begin{proposition}\label{prop::online policy expected throughput upper bound}
% Suppose the memory capacity is not smaller than the equilibrium memory capacity $M^*_j$, consider the expected upper bound of any online policy $c\in \Pi$ in the stochastic scenario, compared with the throughput of the steady system, we have that
% $$\ex{}{\throughput^{c}(T)}\leq \throughput^*= \lambda_j(l_j^\prime+1)$$
% \end{proposition}
% So for any policy without \textbf{Preemption}, the expectation of average throughput is bounded by $\throughput^*$. Detailed proofs are provided in Appendix~\ref{appendix::proof of online policy expected throughput upper bound}
% \begin{comment}The time cost for processing $K\cdot n_j^*$ prompts of type $j$ of static system is $Kd_0+d_{c}KM^*_j$. Under any online policy $c\in \Pi$, the time cost is $\tilde{N}\cdot d_0+d_{c}KM^*_j$ under stochastic scenario, where $\tilde{N}$ is the number of batch iterations. The second term is the same since this term only corresponds to the total memory of all the processed prompts, which is $KM_j^*$. For the first term, it only corresponds to the number of batch iterations to complete the inference of all the prompts. Since the accumulated total memory usage is $KM_j^*$, and the biggest batch size is $M_j^*$, we can easily derive that $\tilde{N}\geq K$. So we have that
% $$\ex{}{\text{Average Throughput}^{c}(T)}\leq \text{Average Throughput}^*= \lambda_j(l_j^\prime+1) $$
% \end{comment}




\subsection{Multiple-Type Fluid Model}
Real systems handle multiple prompt types with different lengths $l_j, l_j'$, competing for memory $C$. Admitting long prompts may force evicting multiple short prompts---a cascading failure analyzed in Section~\ref{sec:wait}.

With $n_j^*$ prompts of type $j$ in equilibrium, total memory is:
\begin{equation}
\label{eq:memory_multi1}
    M^* = \sum_{j=1}^m n_j^*\left(l_j + \frac{l_j^\prime}{2}\right)
\end{equation}
where $(d_0+d_1M^*)\lambda_j(l_j'+1) = n_j^*$ for all $j$. This yields:
\begin{equation}
\label{eq:memory_multi2}
M^* = \frac{d_0\sum_{j=1}^m \lambda_j (l_j^{\prime}+1)\left(l_j + \tfrac{l_j^{\prime}}{2} \right)}{1-d_1\sum_{j=1}^m \lambda_j (l_j^{\prime}+1)\left(l_j + \tfrac{l_j^{\prime}}{2} \right) } 
\end{equation}
The processing time per iteration in equilibrium is:
\begin{equation}    \label{eq:time_fluid}
    \begin{aligned}
            &\text{Processing Time}^* \\&= d_0 + d_1 M^*  \\
                  &= \frac{d_0}{1 - d_1 \sum_{j=1}^m \lambda_j (l_j' + 1)(l_j + \tfrac{1} {2}l_j')},
    \end{aligned}
\end{equation}
as well as the average throughput:
\begin{equation}
    \label{throughput_fluid_multiple}
    \begin{aligned}
        \throughput^*
        &= \sum_{j=1}^m \frac{n_j^* \cdot l_j^{\prime}}{(l_j^{\prime}+1) \cdot (d_0 + d_1 M^*)} \\
        &= \sum_{j=1}^m \lambda_j l_j^{\prime}
    \end{aligned}
\end{equation}
% \begin{equation}
%             \throughput^* 
%         = \frac{1}{d_0}\left(\sum_{j=1}^n n_j^*\right) \cdot \left(1 - \sum_{j=1}^m \lambda_j (l_j^{\prime}+1)\left(l_j + \frac{l_j^{\prime}}{2} \right)\right) 
%         = \sum_{j=1}^m \lambda_j(l_j^{\prime}+1)
% \end{equation}
The equilibrium throughput \(\throughput^*\) benchmarks all non-predictive policies in \(\Pi\).

\begin{proposition}\label{prop::online policy expected throughput upper bound, multiple type}
Suppose that the memory capacity satisfies \(C \geq M^*\), where \(M^*\) denotes the equilibrium memory given in \eqref{eq:memory_multi2}. Then, for any online policy \(\pi \in \Pi\), the expected throughput satisfies
\[
\mathbb{E}[\throughput^{(T,\pi)}] \leq \throughput^*,
\]
where \(\throughput^*\) is given in \eqref{throughput_fluid_multiple}.
\end{proposition}
This quantifies the theoretical limit and reveals a design principle: policies maintaining balanced prompt distribution across stages can approach $\throughput^*$. Our WAIT algorithm (Section~\ref{sec:wait}) uses thresholds to keep the system near equilibrium, preventing cascading failures that plague FCFS (Example~\ref{ex:fcfs_cascade}).




\begin{comment}Similar to the discussion in the single-type fluid model, when the available memory is larger than the equilibrium memory \( M^* \), we can form a batch with a larger batch size, resulting in \(\text{Arrival} < \text{Completion}\). For example, consider the type \( j \) with the smallest value of \((l_j^\prime+1)\left(l_j+\frac{l_j^\prime}{2}\right)\). Without loss of generality, we set \( j=1 \) and the number of prompts of each stage of type $j=1$ is that $\frac{n_1^*}{l_1^\prime+1}+1$. And the memory \( M^+ \) is defined as
$$M^+ = M^* + (l_1^\prime+1)(l_1+\frac{l_1^\prime}{2})$$
So the arrival and completion ratio is that:
\begin{align*}
\frac{ \text{Arrival}}{ \text{Completion} }=& 
\frac{\left(d_0+d_1M^*+d_1(l_1^\prime+1)\left( l_1 + \frac{l_1^\prime}{2} \right)\right)\sum_{j=1}^m \lambda_j}{ \sum_{j=1}^m  \frac{n_j^*}{l_j^\prime+1} +1}\\
=&\frac{ \sum_{j=1}^m  \frac{n_j^*}{l_j^\prime+1} +d_1(l_1^\prime+1)\left( l_1 + \frac{l_1^\prime}{2} \right)\sum_{j=1}^m\lambda_j }{ \sum_{j=1}^m  \frac{n_j^*}{l_j^\prime+1} +1}<1
\end{align*}
The last inequality comes from that: 
$$ d_1(l_1^\prime+1)\left(l_1+\frac{l_1^\prime}{2}\right)\sum_{j=1}^m \lambda_j  \leq d_1\sum_{j=1}^m \lambda_j (l_j^\prime+1)\left(l_j+\frac{l_j^\prime}{2}\right)<1$$
Similarly, when the the available memory is smaller than the equilibrium memory $M^*$, we can form a batch with a smaller batch size, resulting in \(\text{Arrival} > \text{Completion}\). For example, consider the number of prompts of each stage of type \( j=1 \) is that $\frac{n_1^*}{l_1^\prime+1}-1$. And the memory $M^-$ is given by
$$M^- = M^* - (l_1^\prime+1)(l_1+\frac{l_1^\prime}{2})$$
So the arrival and completion ratio is that:
\begin{align*}
\frac{ \text{Arrival}}{ \text{Completion} }=& 
\frac{\left(d_0+d_1M^*-d_1(l_1^\prime+1)\left( l_1 + \frac{l_1^\prime}{2} \right)\right)\sum_{j=1}^m \lambda_j}{ \sum_{j=1}^m  \frac{n_j^*}{l_j^\prime+1} -1}\\
=&\frac{ \sum_{j=1}^m  \frac{n_j^*}{l_j^\prime+1} -d_1(l_1^\prime+1)\left( l_1 + \frac{l_1^\prime}{2} \right)\sum_{j=1}^m\lambda_j }{ \sum_{j=1}^m  \frac{n_j^*}{l_j^\prime+1} -1}>1
\end{align*}
\end{comment}

% with:
% \[
% n_j^* = \lambda_j (1 + l_j') \quad \forall j.
% \]
% Thus:
% \[
% M^* = \sum_{j=1}^m \lambda_j (1 + l_j') \left( l_j + \frac{l_j' + 1}{2} \right)= C,
% \]
% where the last equality holds since the system dynamics are in equilibrium. Aggregating the results above, we get:
% \[
% \throughput^* = \frac{\sum_{j=1}^mn_j^*}{\sum_{j=1}^m\brk{d_0+d_1n_j^*\prn{l_j+\frac{l_j'}{2}}}}.
% \]

% Now we step We refine the fluid model to incorporate this, defining:
% \begin{itemize}
%     \item \(n_{j, \text{prefill}}^*\): steady-state number of type-\(j\) prompts in the prefill stage per batch.
%     \item \(n_{j, \text{decode}}^*\): steady-state number of type-\(j\) prompts in decode stages per batch.
% \end{itemize}

% The \textit{iteration time} \(\tau^*\), the duration to process one batch, is:
% \[
% \tau^* = \max \left\{ d_f \cdot \max_j \{ l_j \}, \; d_1 \cdot \sum_{j=1}^m n_{j, \text{decode}}^* \left( l_j + \frac{l_j' + 1}{2} \right) \right\},
% \]
% where \(d_f\) and \(d_1\) are prefill and decode processing times per token, and \(\frac{l_j' + 1}{2}\) approximates the average decode-stage memory.

% \textbf{Equilibrium Conditions}: The batch processing rate matches arrivals:
% \[
% \frac{n_{j, \text{prefill}}^*}{\tau^*} = \lambda_j \quad \text{and} \quad \frac{n_{j, \text{decode}}^*}{\tau^*} = \lambda_j l_j' \quad \forall j,
% \]
% reflecting one prefill and \(l_j'\) decode iterations per prompt. Memory must satisfy:
% \[
% \sum_{j=1}^m \left( n_{j, \text{prefill}}^* l_j + n_{j, \text{decode}}^* \left( l_j + \frac{l_j' + 1}{2} \right) \right) \leq C,
% \]
% where \(C\) is the GPU memory capacity.

% \textbf{Throughput}: The equilibrium throughput remains:
% \[
% \throughput^* = \sum_{j=1}^m \lambda_j l_j',
% \]
% consistent with token generation rates.

% This refinement aligns the model with practical constraints, assuming a steady batch composition in equilibrium.

% \subsection{Equilibrium Throughput as a Benchmark}


% \arc{Add details on : 1. overloaded system and underloaded system 2. Lemma on upper bound of fluid throughput over all online non-preemptive policy $\Pi$. 3. Check the calculation.}